\paragraph{The group-ring \(L\)-value system for a ramified pair}
Fix a ramified pair \(p=(\nu,u)\) and a sign \(\delta\).
The exponent set \(E_Y\) is \(\{u\}\) for \(u\ge2\) and \(\{0,1\}\) for \(u=1\).
Put \(j_0=0,1\) in the two cases. The complex measure \(F\) is constructed from class sums.

For \(u\ge2\), use the lift from \(U(1)\) to \(U(2u)\). Fix \(Q'\), and take points \(g_a\)
obtained by combining representatives for the source classes at the assigned principal level with
lifts of residue-field units at \(Q'\). For a character \(\tau\) of modulus \(Q'\), the lift is the
sum of the kernel values at the \(g_a\) for the trivial character of that modulus, each multiplied
by \(\tau(\det g_a)^{-1}\). Expanding the two arguments of the Gram pairing gives the squared norm
\[
 \sum_{a,b}\tau(\det(g_b/g_a))\,K_{a,b}.
\]
Here \(K_{a,b}\) is the class sum on the target of the Gram pairings of the two untwisted kernels
at fixed class representatives. A uniform scale \(D_p\) makes each kernel integral.
Clearing the orders of the class stabilizers with \(C_{\mathrm{lev}}\) makes
\(C_{\mathrm{lev}}D_p^2 K_{a,b}\) integral at primes prime to \(\mathrm NQ'\).
Dividing by \((\mathrm NQ')^{2u-1}\) preserves this integrality. These are the coefficients
of \(F\). The sharp rank-one inner product formula gives nonvanishing for \(\tau=1\)
and, by cancellation of the fixed factors, \(F(\xi)/F(1)=\mathcal L_{\nu,u}(\xi;1)\).
When the modulus is enlarged at \(v\), the change removes the factors at \(v,\bar v\)
from the Euler product in the convergence range \(u\ge2\).
\leanlinks{\href{\upstream/grh/src/New/OneLift/HeckeValues/AlgebraicPartIntegrality.lean\#L280}{The double class sum and integrality of the measure}
  \leanlinksep \href{\upstream/grh/src/New/Congruences/Selmer/ClassSums.lean\#L92}{Euler factors and the ratio in the convergence range}}

For \(u=1\), let \(\Lambda(\tau)\) be the determinant period on \(U(2)\) for
\(\nu\tau\upsilon\), where \(\upsilon\) is the fixed splitting character.
At points \(g\) obtained from standard classes and the residue group of the level,
twisting the kernel introduces the factor \(\tau(\det g)\).
The weight polynomial is homogeneous, and at any prime in question a rational translate
makes \(g\) integral. The scale \(D'_p\) for that homogeneous degree therefore makes the
kernel integral, and \(C_{\mathrm{lev}}D'_p\Lambda(\tau)\) is a character sum with integral
coefficients after clearing the class-sum weights.
Take its squared absolute value for \(F\). This is again a measure, by multiplying the
character sum with the sum obtained by conjugating the coefficients and inverting the
evaluation points. The period identity at \(1\) equates a fixed nonzero multiple of
\(|\Lambda(\xi)|^2\) to the product of the norms of the assigned and auxiliary places
times \(|L(1,(\nu\xi)^{-1})|^2\) with the indicated Euler factors omitted.
Comparison with \(\xi=1\) gives \(F(\xi)=F(1)\mathrm NQ'\,\mathcal L_{\nu,1}(\xi;1)^2\).
Adding a place multiplies the period by \((1-X)(1-qX)\), with \(q=\mathrm Nv\) and \(X^{-1}\)
the normalized rank-two uniformizer value. Multiplication by the complex conjugate gives the
Euler factors with exponents \(0,1\) and both signs.
\leanlinks{\href{\upstream/grh/src/New/Congruences/Selmer/ClassSums.lean\#L169}{The period at \(1\), integrality of the measure and the Euler factors}}

In both cases write \(F(\tau)=\sum_j\alpha_j\tau(u_j)\) at modulus \(Q'\).
By Hilbert 90 at the finite components, each \(u_j\) has the same finite part as
\(y_j/\bar y_j\) for a finite idele \(y_j\). Characters of weight zero ignore the infinite components.
Use the finite class of \(\operatorname{divStar}(y_j^{-\delta})\) in place of the point \(u_j\).
Then
\[
 \Theta_{Q'}=\sum_j\iota(\alpha_j)\,[\operatorname{divStar}(y_j^{-\delta})].
\]
Here \(\operatorname{divStar}(y)\) denotes the class of \(y/\bar y\).
The coefficients are integral since \(Q'\) is auxiliary and \(\ell\nmid\mathrm NQ'\).
Evaluation on the Galois character of \(\tau^{-1}\) gives
\(\iota(\tau(u_j)^\delta)\) at each point, by the Artin map.
Thus \(\psi_\tau(\Theta_{Q'})=\iota(F(\tau^\delta))\), and
\(1(\Theta_\varnothing)=\iota(F(1))\ne0\).
\leanlinks{\href{\upstream/grh/src/New/Congruences/Selmer/LValueSystems.lean\#L155}{Passage to the finite class group and character evaluation}}

The character \(\varrho_Y\) is the avatar of the base change of \(\nu^{-\delta}\).
For the norm relation at \(v\), realize a character of the smaller modulus as the avatar of an
inverse tame character, unramified at \(v,\bar v\).
Then \(\varrho_Y(\phi_v)\psi_\tau(\phi_v)\) is the transported \(\delta\)-th power of the
uniformizer value at \(v\) of the base change of \(\nu\tau^\delta\).
The product with both signs makes the exponent \(\delta\) immaterial to the Euler factor.
This gives the norm relation of the system. For a genuine twist \(\xi\),
\(F(\xi^\delta)=F(1)(\mathrm NQ')^{j_0}\mathcal L_{\nu,u}(\xi;\delta)^{j_0+1}\).
The modulus norm is a unit, so only \(j_0+1\) remains as the exponent of the ratio
in the norm of the evaluated group-ring element.

At \(w_0\), the element \(g_0\) is the Artin image of a unit lifting a generator of the
residue-field unit group. The tame character of \(\nu\) is nontrivial there, so
\(\varrho_Y(g_0)\ne1\) and \(\varrho_Y(g_0)^{\mathrm Nw_0-1}=1\).
For a nontrivial root of unity \(\zeta\) whose order \(d\) divides \(n\),
the identity \(d=\prod_{i=1}^{d-1}(1-\zeta^i)\) shows that
\(v_\ell(\zeta-1)\le v_\ell(n)\), since all the other factors are integral.
As \(\#E_Y\le2\), this proves
\(b_Y-v_\ell(1(\Theta_\varnothing))\le4v_\ell(\mathrm Nw_0-1)\).
\leanlinks{\href{\upstream/grh/src/New/Congruences/Selmer/LValueSystems.lean\#L236}{The Euler factors and the reference valuation at \(g_0\)}}
